\chapter{Abstract}

\begin{english}
This paper compares the performance of WebAssembly and JavaScript.
The comparison is made with a clientside backtester that supports trading strategies in both languages.
A backtester is a software that allows the testing of trading strategies based on historical price data.
Strategies are algorithms that automatically make buying and selling decisions in order to make profits without active interaction.
The backtester is a suitable program for comparing the two programming languages, because the backtesting process is very computationally intensive and thus the performance differences are clearly visible.
\newline

The backtester is implemented as a web application in JavaScript that runs in the browser.
To compare the languages, the most computationally intensive part of backtesting, the strategies, is outsourced and implemented in WebAssembly and JavaScript.
To be able to compare the execution times at different computational loads, three different trading strategies with increasing complexity are implemented.
The developed strategies are a moving average crossover strategy, a breakout strategy and a mean reversion strategy with Bollinger bands.
Because of these strategies, the performance and variance of JavaScript and WebAssembly can be compared.
\newline

Overall, it can be seen that the performance of WebAssembly and JavaScript depends on the browser used and the computational intensity of the strategy.
However, WebAssembly can offer a significant advantage in execution for computationally intensive tasks.
Testing the strategies with different data sets shows that WebAssembly in Firefox is faster than JavaScript for computationally intensive strategies.
Furthermore, the variance in the average execution times for WebAssembly in Firefox is lower than that for JavaScript.
In Chrome, on the other hand, JavaScript is faster than WebAssembly for every strategy tested.
The variance in Chrome is lower for WebAssembly than for JavaScript.
\end{english}

